\chapter*{Glossaire}
\addcontentsline{toc}{chapter}{Glossaire}

% ----------------------------------------------------------------
% 📌 À quoi sert cette section ?
% ----------------------------------------------------------------
% Le glossaire permet de regrouper, en fin de document, toutes les définitions de termes techniques ou spécifiques utilisés dans le projet.
% Cela permet d’aider les lecteurs non experts à mieux comprendre certains concepts.

% ----------------------------------------------------------------
% 🧭 Comment utiliser le glossaire dans LaTeX ?
% ----------------------------------------------------------------
% ➤ Déclarez un terme dans le fichier glossaire avec \newglossaryentry
% ➤ Utilisez le terme dans le texte avec \gls{nom}
%
% Exemple dans un fichier .tex :
% \gls{svg} permet de créer des illustrations vectorielles dans le navigateur.
%
% Le mot "svg" apparaîtra en texte clair, et sa définition complète figurera ici, dans le glossaire.

% ----------------------------------------------------------------
% 🧾 Exemple de définition
% ----------------------------------------------------------------

% Déclaration d'un terme dans ce fichier :
\newglossaryentry{svg}{
    name={SVG},
    description={Scalable Vector Graphics. Format d'image vectorielle basé sur XML, utilisé pour le web}
}

% Déclaration d'un second terme :
\newglossaryentry{dompurify}{
    name={DOMPurify},
    description={Bibliothèque JavaScript permettant de nettoyer du code HTML/SVG pour éviter les failles XSS}
}

% ----------------------------------------------------------------
% ✍️ À vous de jouer :
% ----------------------------------------------------------------
% ➤ Dupliquez les blocs \newglossaryentry pour chaque mot à définir.
% ➤ Utilisez \gls{mot} dans vos fichiers .tex pour faire référence à ces définitions.

